\documentclass[10pt,a4paper]{amsart}
\usepackage[margin=19mm]{geometry}
\usepackage{amsmath,amssymb,mathtools}
\usepackage[scr=rsfs,cal=euler]{mathalfa}
\usepackage{xfrac}
\usepackage[unicode,hidelinks,bookmarks=false]{hyperref}
\newcommand{\ie}{i.e.\ }
\newcommand{\cf}{cf.\ }
\newcommand{\resp}{respectively\ }
\newcommand{\Subsection}{\subsection}
\providecommand{\nobreakdash}{\nobreak\hbox{-}}
\setlength{\emergencystretch}{2em}
\multlinegap0pt
\usepackage{amssymb}
\usepackage{enumerate}
\newtheorem{lemma}{Lemma}
\title{Universal vector bundles, push-forward formul{\ae} and positivity of characteristic forms}
\author{Filippo Fagioli}
\date{Source: arXiv:2210.11157v1 (CC BY 4.0)}
\begin{document}
\maketitle

\noindent\emph{Source and licence.} Universal vector bundles, push-forward formul{\ae} and positivity of characteristic forms. Version arXiv:2210.11157v1 (CC BY 4.0), distributed by arXiv under the Creative Commons Attribution 4.0 International licence, http://creativecommons.org/licenses/by/4.0/, as recorded in the arXiv licence metadata for this exact version and shown on the abstract page https://arxiv.org/abs/2210.11157v1. Full licence notice: \texttt{licenses/arxiv-2210.11157-CC-BY-4.0.txt}.\par\smallskip

\noindent\textit{Editorial scope.} Counted unit: the article's lemma lem:blocksgeneral (source0.tex lines 611--613). The article's complete original proof of it is delivered with it (source0.tex lines 614--647, one \texttt{proof} environment of 34 lines). No mathematics is written, repaired or re-derived: the statement and the proof are the article's own lines, in the article's own notation, and no retained line is substituted or re-flowed.  No further source-local context is needed: every cross-reference of every retained line resolves inside the two delivered spans.  Nothing was omitted from any retained span, so the package carries no omission record.  No retained line contains a citation, so the wrapper carries no bibliography and no bibliography entry.  No retained line contains a figure environment, an image-inclusion command or drawing code, so no image is delivered and the package declares no asset dependency.  Editorial formatting only, in addition: an AMS \texttt{amsart} preamble that restates the article's own declarations and macros used by the retained lines, namely the environments \texttt{lemma} on separate counters, together with the standard packages those lines need.  The retained environments print in this document as Lemma 1 in order of appearance, whereas the article prints the same items with the numbering of its own full document.  The complete native source of the article as distributed in the arXiv e-print for this exact version is bundled unchanged as \texttt{sources/132349/source0.tex}.
\begin{lemma}\label{lem:blocksgeneral}
	The section $ \theta_{(\ell,l)} $ is well defined, \textsl{i.e.} it does not depend upon the choice of a particular representative $ \mathbf v= (v_1,\dots,v_r) $ for $ \mathbf{f} $.
\end{lemma}

\begin{proof}
	Take a local normal frame $ (e_1,\dots,e_r) $ at $ x $ such that $ \mathbf{e} := \bigl(e_1(x),\dots,e_r(x)\bigr) $ and $ \mathbf{v} $ identify the same flag $ \mathbf{f} $.
	Consequently, we know that if an index $ \lambda $ is such that
	\[
	r-\rho_l < \lambda \le r-\rho_{\ell}
	\]
	then we can write
	\begin{equation*}
		v_\lambda = \sum_{r-\rho_l < \alpha \le r-\rho_{\ell}} a_{\alpha \lambda} e_{\alpha}(x),
	\end{equation*}
	where the coefficients $ (a_{pq}) $ are the entries of the change of coordinates matrix
	\begin{equation*}
		\begin{pmatrix}
			A_{11} & 0 & 0 \\
			0 & \ddots & 0 \\
			0 & 0 & A_{mm}
		\end{pmatrix}.
	\end{equation*}
	Notice that the diagonal block $ A_{jj} $ is again a unitary matrix of size $ \rho_{m-j+1} - \rho_{m-j} $.
	
	To simplify the notation in what follows, we omit the dependence on $ x $ (resp. $ (x,\mathbf{f}) $) of $\mathbf{e}$ (resp. $  \pi_{\rho}^*\Theta(E,h)_{(x,\mathbf{f})} $).
	We thus get the chain of equalities
	\begin{align*}
		& \frac{i}{2\pi} \sum_{\lambda,\mu} \bigl\langle \pi_{\rho}^*\Theta(E,h)\cdot {v_\lambda}, {v_\mu} \bigr\rangle_h \otimes v_{\lambda}^{\vee} \otimes v_{\mu} \\
		&= \frac{i}{2\pi} \sum_{\lambda,\mu} {\left\langle {\pi_{\rho}^*\Theta(E,h)}\cdot \left( \sum_{\alpha} a_{\alpha \lambda} e_\alpha \right), \sum_{\beta} a_{\beta \mu} e_\beta \right\rangle_h} \otimes \left( \sum_{\tilde{\alpha}} a_{\tilde{\alpha} \lambda} e_{\tilde{\alpha}} \right)^{\vee} \otimes \sum_{\tilde{\beta}} a_{\tilde{\beta} \mu} e_{\tilde{\beta}} \\
		&= \frac{i}{2\pi} \sum_{\alpha,\tilde{\alpha}}\underbrace{ \sum_{\lambda} a_{\alpha \lambda} \overline{a_{\tilde{\alpha} \lambda}}}_{=\delta_{\alpha\tilde{\alpha}}} \sum_{\tilde{\beta},\beta}\underbrace{ \sum_{\mu} a_{\tilde{\beta} \mu} \overline{a_{\beta \mu}}}_{=\delta_{\tilde{\beta}\beta}} \,\bigl\langle {\pi_{\rho}^*\Theta(E,h)}\cdot {e_\alpha}, {e_\beta} \bigr\rangle_h \otimes e_{\tilde{\alpha}}^{\vee} \otimes e_{\tilde{\beta}} \\
		&= \frac{i}{2\pi} \sum_{\alpha,\beta}\bigl\langle {\pi_{\rho}^*\Theta(E,h)}\cdot {e_\alpha}, {e_\beta} \bigr\rangle_h \otimes e_{\alpha}^{\vee} \otimes e_{\beta},
	\end{align*}
	where the indices of the summations are such that
	\[
	\lambda,\mu,\alpha,\beta,\tilde{\alpha},\tilde{\beta} = r-\rho_l + 1, \dots, r-\rho_{\ell} .
	\]
	Hence the lemma follows.
\end{proof}

\end{document}
